\section{A Counterexample to NP03 Proposition 16}
\label{app:np03}

This appendix uses a binary covariant constructor $\ftree$ and the child
alphabet $\{0,1\}$. NP03's binary child indices $1,2$ become $0,1$ here.
A separate kernel-checked formalization of this counterexample, not part
of the artifact (\cref{app:lean}), agrees with the argument below.

\subsection{Cap Automata and the Back-Translation}

A \emph{cap automaton} consists of a finite automaton and an additional relation
of \emph{P-edges}, written $q'\pedge q$. Write
$q\xrightarrow{w}q'$ for a run of the underlying finite automaton.
For one initial state $q_0$ and final-state set $F$, its language is
\[
\Lang(P)=
 \{w\mid q_0\xrightarrow{w}q,\ q\in F\}
 \;\cup\;
 \{\pi u\mid q_0\xrightarrow{\pi}q\xrightarrow{\rho}q',\
              q'\pedge q,\ u\in\pr(\rho^*)\}.
\]
Here $\rho^*$ is the set of finite powers of $\rho$, and $\pr$ takes all
prefixes (\cref{sec:words}). Thus a P-edge accepts prefixes of repetitions of a word read
between its target and its source. We allow $\rho=\eps$, as in NP03's
Definition~1; the example below has no cap at all and also works if roots
are required to be nonempty.
A cap automaton is \emph{complete} if every word has a run from an initial
state, whether or not that run accepts. NP03's Definition~3 calls it
\emph{restricted} when the following four conditions hold.
\begin{enumerate}[leftmargin=*,nosep]
\item \emph{Strictly epsilon free:} there is one initial state and no
epsilon transition.
\item \emph{Gap universal:} if $q\xrightarrow{i}q'$ with $q\notin F$
and $q'\in F$, then every word has a run from $q'$ to a final state.
This condition uses ordinary finite-automaton acceptance.
\item \emph{Strictly cap:} if
$q\xrightarrow{\rho}q'\pedge r$ with $\rho\ne\eps$, then $q\in F$.
\item \emph{Shuffled:} if
$q_0\xrightarrow{\pi}q$, $q_0\xrightarrow{\pi}q'$, and
$q'\pedge q$ with $q\ne q'$, then $\pi\tau\in\Lang(P)$ for every word
$\tau$.
\end{enumerate}

NP03's Definition~5 introduces distinct variables $l(q),r(q)$ for each state,
and defines $J(P)=(\judg{\phi_P}{l(q_0)}{r(q_0)})$.
The binary rules needed here are as follows. For each choice of successors
$q\xrightarrow{0}p_0$ and $q\xrightarrow{1}p_1$, include
\[
 l(q)\le\ftree(l(p_0),l(p_1))\quad(q\in F),
 \qquad
 \ftree(r(p_0),r(p_1))\le r(q).
\]
A transition $q\xrightarrow{i}q'$ with $q\notin F$ and $q'\in F$
contributes $r(q')=\Top$. Finally,
$q\xrightarrow{i}q'\pedge p$ with $q'\ne p$ contributes
\[
 l(q)\le\ftree(r(p),t)\quad(i=0),\qquad
 l(q)\le\ftree(t,r(p))\quad(i=1),\qquad t=\Top,
\]
where $t$ is a fresh filler variable. These last constraints are the
binary form of the notation $l(q)\le i[r(p)]$ in Table~7.
NP03 does not fix the other argument of $i[\cdot]$. Leaving $t$
unconstrained gives an equivalent judgment: it occurs only inside this
upper bound, so raising it to $\Top$ preserves every solution without
changing $x,y$. Conversely, a solution with $t=\Top$ satisfies the
unconstrained version. Thus both versions have the same solutions on
the original variables in each tree domain.

In Table~7, the finality side conditions of the left and top rules are
conveyed by typographic marks. We read the left rule as applying at final
sources and the top rule at final targets, as NP03's proof of
Proposition~16 uses in the auxiliary claim in case~(3) and the
gap-universality argument in case~(3b). Dropping these conditions adds
left bounds at nonfinal sources and top equations at nonfinal targets.
We prove the counterexample for both readings and every
constraint system between them, so the conclusion does not depend on
these marks.

\subsection{The Six-State Automaton}

Let the states be $q_0,A,B,C,C_2,U$, with initial state $q_0$,
final states $q_0,A,U$, and the single P-edge $C\pedge B$.
The following table lists all letter transitions; an entry denotes the
set of possible successors.
\begin{center}
\begin{tabular}{@{}lcc@{}}
\toprule
State & on $0$ & on $1$\\\midrule
$q_0$ & $\{U\}$ & $\{A,B\}$\\
$A$   & $\{U\}$ & $\{C\}$\\
$B$   & $\{U\}$ & $\{C_2\}$\\
$C$   & $\{U\}$ & $\{C_2\}$\\
$C_2$ & $\{U\}$ & $\{C_2\}$\\
$U$   & $\{U\}$ & $\{U\}$\\\bottomrule
\end{tabular}
\end{center}
Every state has at least one successor on each letter, so $P$ is complete.
It has one initial state and no epsilon transitions. Every transition
from a nonfinal to a final state enters $U$, which accepts every
continuation by ordinary acceptance. The only states with a nonempty
run to the P-edge source $C$ are $A$ and $q_0$, both final. Finally,
$B$ is reachable from $q_0$ only on $1$, and $C$ only on $11$.
No word reaches both endpoints of the P-edge, so shuffledness is
vacuous. This checks all four restrictions.

Every word containing $0$ reaches $U$. On $1^j$, the reachable state
sets are $\{q_0\}$ for $j=0$, $\{A,B\}$ for $j=1$,
$\{C,C_2\}$ for $j=2$, and $\{C_2\}$ for $j\ge3$.
No run from $B$ reaches $C$, so no P-edge contributes to acceptance.
Consequently
\[
 \Lang(P)=\{0,1\}^*\setminus\{1^j\mid j\ge2\},
 \qquad 11\notin\Lang(P).
\]

\subsection{The Generated Constraints and Their Solutions}

Write $l_q=l(q)$ and $r_q=r(q)$, and set $x=l_{q_0}$ and $y=r_{q_0}$.
With the final-state marks, Table~7 gives the following conjunction
$\phi_{\min}$; the grouping records which rule supplies each literal.
\begin{align*}
\text{left:}\quad
 x&\le\ftree(l_U,l_A), & x&\le\ftree(l_U,l_B),\\
 l_A&\le\ftree(l_U,l_C), & l_U&\le\ftree(l_U,l_U);\\[2pt]
\text{right:}\quad
 \ftree(r_U,r_A)&\le y, & \ftree(r_U,r_B)&\le y,\\
 \ftree(r_U,r_C)&\le r_A, & \ftree(r_U,r_{C_2})&\le r_B,\\
 \ftree(r_U,r_{C_2})&\le r_C, & \ftree(r_U,r_{C_2})&\le r_{C_2},\\
 \ftree(r_U,r_U)&\le r_U;\\[2pt]
\text{top:}\quad r_U&=\Top;\\
\text{P-edge:}\quad l_A&\le\ftree(t,r_B), & t&=\Top.
\end{align*}
In particular, the P-edge literal comes from
$A\xrightarrow{1}C\pedge B$.
Without the final-state marks, the system $\phi_{\max}$ adds precisely
\[
 l_B\le\ftree(l_U,l_{C_2}),\quad
 l_C\le\ftree(l_U,l_{C_2}),\quad
 l_{C_2}\le\ftree(l_U,l_{C_2}),\quad r_{C_2}=\Top.
\]
Set every $l_q$ to $\Bot$ and every $r_q$ and $t$ to $\Top$.
Each upper bound then has $\Bot$ on the left, each lower bound has
$\Top$ on the right, and all top equations hold. This is a finite
solution of $\phi_{\max}$, hence of any system $\phi$ with
$\phi_{\min}\subseteq\phi\subseteq\phi_{\max}$.

We next prove $\phi_{\min}\models x\le y$ over arbitrary trees.
Fix a solution $\sigma$ and abbreviate its values by
$X=\sigma(x)$, $Y=\sigma(y)$, $a=\sigma(l_A)$,
$b=\sigma(r_B)$, and $c=\sigma(r_{C_2})$.
The lower bounds give
\[
 \ftree(\Top,c)\le c,\qquad
 \ftree(\Top,c)\le b,\qquad
 \ftree(\Top,b)\le Y.
\]
They imply that every label of $Y$ is either $\Top$, or $\ftree$ at a
path of the form $1^j$. Here is a direct verification.
From $\ftree(\Top,c)\le c$, $c$ is either $\Top$ or
$\ftree(\Top,c')$ with $c\le c'$. In the latter case,
substitution gives $\ftree(\Top,c')\le c'$. Repeating this argument
along any finite path proves the asserted label property for $c$.
That property is preserved by $s\mapsto\ftree(\Top,s)$ and by passing
to a larger tree: at a constructor root the larger tree either becomes
$\Top$ or keeps a top left child and a larger right child.
Induction on path length justifies the latter assertion even for infinite
trees. Applying these two preservation facts to the remaining lower
bounds proves the property for $Y$.

For a tree $s$, let $s|_w$ denote its subtree at $w$, when $w$ belongs
to its domain, and write $s(w)$ for its label there. We check the labels of $X$ and $Y$ at every common
path $w$, which suffices for covariant subtyping. If $Y(w)=\Top$,
the labels compare. Otherwise, $w=1^j$ and $Y(w)=\ftree$.
For $j=0$, the bound $X\le\ftree(\sigma(l_U),a)$ excludes a top root
of $X$. For $j=1$, that bound and
$a\le\ftree(\sigma(l_U),\sigma(l_C))$ exclude a top label of $X$ at $1$.
For $j\ge2$, the existence of $1^j$ in $X$ makes its root and its
child at $1$ constructors. The bounds
\[
 X\le\ftree(\sigma(l_U),a),\qquad
 a\le\ftree(\Top,b),\qquad
 \ftree(\Top,b)\le Y
\]
then force $a$ to be a constructor too: $X|_1\le a$ excludes $\Bot$,
and the upper bound on $a$ excludes $\Top$.
Since $1^j$ also exists in $Y$, its root is a constructor.
Decomposing these three inequalities gives the subtree inequalities
\[
 X|_{11}\le a|_1\le b\le Y|_1.
\]
By transitivity, $X|_{11}\le Y|_1$. Comparing at their common path
$1^{j-2}$ yields
$X(1^j)\le Y(1^{j-1})=\ftree$, because $1^{j-1}$ is a proper prefix
of a path in $Y$. Thus $X(w)\le Y(w)$ in every case, proving $X\le Y$.
Entailment persists when literals are added, so every intermediate
reading $\phi$ is both satisfiable and entails $x\le y$.
The proof covers arbitrary trees and hence also regular and finite trees.

\subsection{Forward Acceptance and the Failed Proof Step}

Write $P_{\mathsf l}(J(P))$ and $P_{\mathsf r}(J(P))$ for the left and
right cap automata obtained from $J(P)$ by NP03's Table~3.
The rules used here can be read directly on these constraints: support
between variables is trivial, as in NP03's Lemma~21.
States are pairs $(u,v)$ of variables or a blank
coordinate, with initial state $(x,y)$, and an accepting sink.
An upper bound $u\le\ftree(u_0,u_1)$ and a lower bound
$\ftree(v_0,v_1)\le v$ permit a transition on $i$ to $(u_i,v_i)$;
either coordinate may instead descend alone, leaving the other blank.
The left automaton makes a state final when its left coordinate has an
upper constructor bound; the right automaton uses a lower constructor
bound on the right coordinate. The bottom, top, and reflexivity rules
send ready states to the accepting sink on every letter. In particular,
a right coordinate constrained to $\Top$ is ready.
The left P-edge rule includes $(u,\square)\pedge(v,u)$, where
$\square$ denotes a coordinate that may be a variable or blank.

The left automaton of $J(P)$ has the transitions and P-edge
\[
 (x,y)\xrightarrow{1}(l_A,r_B)
       \xrightarrow{1}(r_B,\text{blank})
       \pedge(l_A,r_B).
\]
The first transition combines the upper bound using $l_A$ with the lower
bound using $r_B$; the second uses the P-edge literal.
This cap uses only descend steps and a P-edge. Its entry state is final
because $l_A\le\ftree(l_U,l_C)$ is an upper constructor bound.
Hence $11\in\Lang(P_{\mathsf l}(J(P)))$ under either finality convention
for Table~3, whereas $11\notin\Lang(P)$. The inclusion
$\Lang(P_{\mathsf l}(J(P)))\subseteq\Lang(P)$, claimed in case~(3) of
NP03's proof of Proposition~16, therefore fails independently of that
convention, even if cap entry states are required to be final.

To show universality, we additionally count ready states as final.
NP03's proofs of Propositions~12 and~14 use this convention:
Proposition~12 uses bottom, top, and reflexivity readiness to accept
the prefix itself; Proposition~14 uses top-readiness for gap universality
and reflexivity for shuffledness. NP03's own Section~4 example witnesses
the need for this reading of Table~3. For the unary judgment
$\judg{z=\Top\wedge f(z)\le y}{x}{y}$, NP03 states that every path in
$1^+$ is both left safe and right safe, using its child index $1$.
Read literally, Table~3 reaches only $(\text{blank},z)$ on $1$;
that state is not final in either automaton, and there is no cap, so
both reject $1$.

Under this convention, the displayed cap accepts $1\pr(1^*)$, and
$(x,y)$ is final, accepting $\eps$.
For a word containing $0$, let the right coordinate follow any run of
$P$ through the lower constructor bounds. At its first $0$, it reaches
$r_U=\Top$, so this path and all its extensions are accepted.
The right automaton is universal as well: since every state of $P$ has
a successor on each letter, a right-coordinate run exists on every word,
and every $r_q$ has a lower
constructor bound, so every resulting state is final.
These runs remain available under every intermediate reading of Table~7.
They verify directly the conclusion also supplied by NP03's
Propositions~1 and~3 (soundness, Proposition~8; completeness,
Proposition~11):
\[
 \Lang(P_{\mathsf l}(J(P)))=\{0,1\}^*\ne\Lang(P),
 \qquad \Lang(P_{\mathsf r}(J(P)))=\{0,1\}^*.
\]
The convention is needed only for this universality conclusion; acceptance
of $11$ already refutes Proposition~16.

Case~(3d) of NP03's proof writes the cap's entry state as
$(l(q_i),r(q_i))$, assuming that both coordinates follow the same run
of $P$. In this example the entry state is $(l(A),r(B))$.
Both $A$ and $B$ are reached on $1$, but only the run through $A$
can next reach $C$, whose P-edge returns to $B$.
The back-translation combines coordinates of different runs; the original
automaton has no run $B\xrightarrow{1}C$ to close that cap.
Shuffledness does not apply: it tests endpoints reached on the same
word, whereas $B$ and $C$ occur at different cuts here.

\subsection{Consequences for the Published Reductions}

Theorem~4 explicitly invokes Proposition~16 to reduce restricted
cap-automaton universality to entailment. The counterexample invalidates
that construction and therefore the converse reduction underlying
Theorem~1. This is the implication from decidability of NSSE to
decidability of restricted cap-set universality. It does not rule out a
different reduction between the two decision problems.
Corollary~1, asserting equivalence of arities at least two and of finite,
regular, and arbitrary tree classes, also relies on that converse.
Its arity part additionally uses Proposition~17's alphabet extension and
its alphabet restriction, which encodes each letter by two letters of a
smaller alphabet. Theorem~4, Theorem~1, and Corollary~1 themselves are
not refuted; their published proofs are incomplete. We claim no defect
in Proposition~17.

The forward reduction of Theorem~2 first checks satisfiability in the
chosen tree domain, then constructs the two Table~3 automata.
Its soundness and completeness proofs precede the back-translation and
do not invoke Proposition~16. The unary result, Corollary~3, follows
in NP03 from Corollary~2: cap-set membership is expressed by existential
word equations with regular constraints, and universality adds a
universal word quantifier. Although the text introduces this derivation
through Theorem~1, it uses only the forward reduction to cap sets.
For a unary alphabet, word equations and regular constraints become
constraints in Presburger arithmetic on lengths, so this quantified
formulation is decidable.
Neither the converse reduction nor Proposition~17 is needed.
